\chapter*{General Conclusion}
\addcontentsline{toc}{chapter}{General Conclusion}

\markboth{General Conclusion}{General Conclusion}

Riaya was built at Pengguin, a company led by John Jang, for patients in Tunisia who find it hard to book a doctor. Many of them call cabinet after cabinet during working hours, are not at ease with websites and apps, or speak mostly Tunisian Arabic. With Riaya they call one phone number, talk to an AI assistant in Derja, French or English, and receive the rest on WhatsApp. Doctors manage their practice from a web dashboard, and administrators verify doctors and monitor the calls.

\section*{Summary of the Work Completed}

The first chapter presented Pengguin and the Riaya product: its actors (the patient without an account, the doctor, the secretary and the administrator), and the Free and Pro plans. It explained how each plan limit is counted and why phone booking is never limited. It then set out the functional and non-functional requirements. A comparison with Doctolib, the Tunisian directories and classic IVR systems showed that none of them lets a patient book by speaking in Derja, sends urgent requests to several doctors, or follows up on WhatsApp.

The second chapter covered the work method and the models. The project followed Kanban on a Notion board shared by John Jang, the academic supervisor and me. The cards on the board show how requests arrived: most came from John, who was the only one in contact with the client, and research cards such as the comparison of speech models stayed open until they had a written conclusion. Kanban suited a single developer, requests that came at irregular times and frequent feedback from the client better than Scrum would have. The second part modeled the platform with use case, class, sequence, activity and deployment diagrams, and the BPMN booking process.

The third chapter described the architecture and its implementation. The code is split into a Next.js web application, which owns the PostgreSQL database and the business rules, and an Express voice server, which carries the phone audio between Twilio and Gemini Live, sends WhatsApp messages and runs the background jobs. Places named by callers are found with Google Places, and the best-fit search ranks nearby doctors by distance and waiting time before a road-distance check on the best eight. The chapter explained how phone booking works, from latency to urgent requests, callbacks and reviews, along with the decisions to keep Free doctors in the search and not to show prices, and how consultation recordings and the AI assistant work. It then presented the evaluations that decided the place search provider, the road-distance step and the voice model, with how the platform was tested, went through the security measures and their open gaps, and showed the screens of the four interfaces.

\section*{Challenges Encountered and Lessons Learned}

Tunisian Arabic was the hardest part of the project. Few models understand it, and the sentences that mix French words into Derja gave the most errors. We compared realtime models and chains of separate models for listening, thinking and speaking, using a small program that plays the same seven sentences against the real prompt and tools. The tests showed a trade-off between speed and quality: the models with the best understanding and the most natural voice were often the slowest, and in a chain the delays add up. The fastest chain we built still waited about 6.8 seconds before answering. Gemini Live gave the most natural Tunisian dialect and answered fast enough for a phone call, so it became the final choice. Its own transcript of the caller was very inaccurate, though, so we added a separate transcriber for the caller's words. These tests used synthetic voices, and real patients will be the true test of the choice.

Geocoding taught a similar lesson. Callers name hospitals, schools, stations and streets in Derja or French, and the search is only as good as the position it gets. We checked several providers on real places whose position we knew, and Google Places found more of them and placed them more accurately. The road-distance step was measured in the same spirit: over 78 comparisons it changed the best doctor in about a third of the cases, but never saved a patient five minutes or more of driving, so it stays as a light polish.

Urgent requests and callbacks needed careful design. The AI tells a caller with life-threatening signs to call the SAMU first, then sends up to three urgent requests that cancel each other once a doctor accepts. Delayed jobs cancel unanswered requests after ten minutes and call the patient back, and a refusal leads to a new call that offers other doctors. Each of these flows has several edge cases (a blocked number, a patient who does not pick up, a refused urgent request). The activity diagrams of \chapref{chap:management} follow each of them.

Some work is still unfinished. The code has no automated test suite yet, so a regression is caught only when the feature is tested again, manually or with AI tools. Some routes of the voice server, including the incoming-call webhook and the live-call WebSockets, do not yet check a signature or a shared secret.

\section*{Future Perspectives}

The first task is automated testing, starting with the best-fit search, the plan limits and the booking flow, which are the parts where a silent regression would hurt patients or doctors most. The second is securing the voice server: verifying Twilio's signature on the incoming-call webhook and requiring the shared secret or a session on the WhatsApp route and the live WebSockets. The third concerns files. Medical documents and recordings are currently reached through public CDN links with random names; moving them to a private bucket and handing out short-lived signed download links after the ownership check would close that gap. The transcription calls should also receive the same no-retention options as the chat calls.

The Free-plan limits were set before real usage existed. Once enough doctors use the platform, these limits should be tuned against what doctors actually do and what each feature costs.

Further ahead, the design could serve other countries. The voice server already keeps several model backends behind one setting, and the prompts, the place search and the payment provider are the parts that would change. Teleconsultation by video would let a doctor see some patients without a visit, and richer analytics, such as booking and no-show trends, would help doctors plan their week.
